\subsection{A minimal identifying experiment and precision planning}\label{sec:minimal}
The smallest targeted experiment tests the shared-phase observation model, rather than every extension at once. Independently establish a coarse delay branch, record calibrated carrier offsets at two distinct frequencies, and scan receptive gain at one frequency with matched content and power. Fit $\gamma,\alpha,\tau$ using the two carrier offsets and the receptive offset. Predict an independently measured onset latency after separately accounting for processing. A discrepancy beyond the frozen uncertainty rejects the shared transfer/processing assumptions or the observation channel. A single band supplies an equivalence class, so there is no uniquely defined single-band delay estimator whose systematic bias must be demonstrated to confirm the ridge.

For fixed known amplitude and concentration, one probe away from the stationary points of the gate has nonzero local receptive information. This mathematical minimum is a poor practical design: an unknown gate amplitude can confound it. Use several phases spanning both sides of the center, reserve phases for held-out prediction, and estimate nuisance parameters jointly. The supplied illustration uses 48 scan phases, 32 independent phase trials per frequency, and two frequencies; these are generator choices, not established biological minima. Computing a two-band offset fit is constant-size linear algebra; a nonlinear fit with $N$ probe observations costs $O(N)$ per residual/Jacobian evaluation for a fixed number of parameters. No algorithmic complexity advantage is claimed.

\paragraph{Timing precision is not an electrode count.} For independent latency estimates with standard deviation $\sigma_\tau$, a normal-approximation 95\% mean-interval half-width $e$ requires $N\gtrsim(1.96\sigma_\tau/e)^2$. With illustrative $\sigma_\tau=3$ ms and $e=1$ ms, this is 35 independent estimates. Repeated pulses and nearby electrodes need not be independent; subject-level uncertainty requires subjects or a justified hierarchical model. At least one valid stimulus site and a recorded target are necessary for a directed timing test, and both directional interventions are needed for reciprocity. There is no universal number of electrodes per hemisphere that ensures delay precision. Spatial coverage, artifacts, sampling, pulse variability and indirect routes determine eligibility and effective information. Estimate these from pilot waveforms before freezing sample size.

\paragraph{Factorial contrast precision.} For a prespecified scalar receiver contrast, equal independent cell sizes $n$ and within-cell variance $\sigma^2$ give variance $4\sigma^2/n$ for the four-cell difference. To detect magnitude $D$ with power $1-\beta$ using $K$ prespecified two-sided contrasts and a Bonferroni family bound $\alpha_0$,
\begin{equation}
n\gtrsim\frac{4\sigma^2}{D^2}
\left[z_{1-\alpha_0/(2K)}+z_{1-\beta}\right]^2,
\label{eq:factorial-power}
\end{equation}
where $z_p$ is the standard normal quantile. For $K=4$, $\alpha_0=0.05$, power 0.8, and $D/\sigma=0.5$, this gives 179 trials per cell. It is a planning illustration for independent Gaussian contrasts, not a recommendation for Allen recordings. Ordinary stimulus classes are not automatically the four independent pathway interventions required by the factorial theorem.

The number of basis pairs is $r_1r_2$, but rank detection also depends on the smallest nonzero singular value and uncertainty in maps, whitening, and contrasts. A single nonzero scalar contrast does not identify rank. If a simultaneous calibration supplies $\|\widehat{QC_B}-QC_B\|_2\le e_C$ with its declared confidence, singular-value perturbation bounds give $\sigma_j(\widehat{QC_B})>e_C$ as evidence that the true $j$th singular value is positive on that event. An observed value below $e_C$ is unresolved, not evidence of an exact null. Pilot cluster resampling and held-out null controls must determine $e_C$ and the achievable effect bound.

\paragraph{Decisive component tests.} Test identifying assumptions by predicting independent timing; test content delivery with an independently calibrated sender contrast at matched phase; test returns by changing the reverse pathway while checking that the initial forward increment is unchanged; and test learning by training one content subspace and evaluating an orthogonal held-out direction. Nonzero projection without threshold crossing and two nonzero maps without an item return are predicted failure cases, not falsifications. Failure to reach statistical significance alone does not reject a model: an exclusion interval against a prespecified meaningful effect is required.
